% Compile with: pdflatex supplement_27_1.tex
% Requires: amsmath, amssymb, amsthm, mathtools, xcolor, mdframed,
%           enumitem, hyperref, pagecolor, microtype
\documentclass[11pt]{article}
\usepackage[margin=1in]{geometry}
\usepackage{amsmath, amssymb, amsthm}
\usepackage{mathtools}
\usepackage{xcolor}
\usepackage{mdframed}
\usepackage{enumitem}
\usepackage{hyperref}
\usepackage{pagecolor}
\usepackage{microtype}

% ── Colour palette ──────────────────────────────────────────
\definecolor{cream}{HTML}{FFFDF5}
\definecolor{defblue}{HTML}{1A4A7A}
\definecolor{thmgreen}{HTML}{1A5C3A}
\definecolor{warnAmber}{HTML}{7A4A00}
\definecolor{dangerRed}{HTML}{7A1A1A}
\definecolor{boxbluebg}{HTML}{EBF3FB}
\definecolor{boxgreenbg}{HTML}{EBF7EF}
\definecolor{boxamberbg}{HTML}{FDF5E6}
\definecolor{boxredbg}{HTML}{FBEBEB}
\pagecolor{cream}

% ── Theorem styles ──────────────────────────────────────────
\newtheoremstyle{defstyle}  {8pt}{8pt}{\normalfont}{0pt}{\bfseries\color{defblue}} {.}{0.5em}{}
\newtheoremstyle{thmstyle}  {8pt}{8pt}{\normalfont}{0pt}{\bfseries\color{thmgreen}}{.}{0.5em}{}
\newtheoremstyle{notstyle}  {8pt}{8pt}{\normalfont}{0pt}{\bfseries\color{warnAmber}}{.}{0.5em}{}
\newtheoremstyle{cautionstyle}{8pt}{8pt}{\normalfont}{0pt}{\bfseries\color{dangerRed}}{.}{0.5em}{}

\theoremstyle{defstyle}
\newtheorem{definition}{Definition}[section]

\theoremstyle{thmstyle}
\newtheorem{theorem}{Theorem}[section]
\newtheorem{corollary}[theorem]{Corollary}

\theoremstyle{notstyle}
\newtheorem*{remark}{Remark}
\newtheorem*{example}{Example}

\theoremstyle{cautionstyle}
\newtheorem*{caution}{Caution}

% ── Framed boxes ────────────────────────────────────────────
\newmdenv[
  backgroundcolor=boxbluebg, linecolor=defblue, linewidth=2pt,
  topline=false, rightline=false, bottomline=false,
  innerleftmargin=10pt, innerrightmargin=10pt,
  innertopmargin=6pt, innerbottommargin=6pt,
  skipabove=8pt, skipbelow=8pt
]{defbox}

\newmdenv[
  backgroundcolor=boxgreenbg, linecolor=thmgreen, linewidth=2pt,
  topline=false, rightline=false, bottomline=false,
  innerleftmargin=10pt, innerrightmargin=10pt,
  innertopmargin=6pt, innerbottommargin=6pt,
  skipabove=8pt, skipbelow=8pt
]{thmbox}

\newmdenv[
  backgroundcolor=boxamberbg, linecolor=warnAmber, linewidth=2pt,
  topline=false, rightline=false, bottomline=false,
  innerleftmargin=10pt, innerrightmargin=10pt,
  innertopmargin=6pt, innerbottommargin=6pt,
  skipabove=8pt, skipbelow=8pt
]{notebox}

\newmdenv[
  backgroundcolor=boxredbg, linecolor=dangerRed, linewidth=2pt,
  topline=false, rightline=false, bottomline=false,
  innerleftmargin=10pt, innerrightmargin=10pt,
  innertopmargin=6pt, innerbottommargin=6pt,
  skipabove=8pt, skipbelow=8pt
]{cautionbox}

% ── Macros ──────────────────────────────────────────────────
\newcommand{\Z}{\mathbb{Z}}
\newcommand{\N}{\mathbb{N}}
\newcommand{\R}{\mathbb{R}}
\newcommand{\id}{\mathrm{id}}

\begin{document}

% ── Title ───────────────────────────────────────────────────
\begin{center}
  {\Large\bfseries Supplement 27.1 --- $|\N| = |E|$}\\[4pt]
  {\normalsize CS 173: Discrete Structures \quad$\cdot$\quad Functions \& Cardinality}\\[2pt]
  {\small\color{gray} University of Illinois at Urbana-Champaign}
\end{center}
\vspace{1em}\hrule\vspace{1.5em}

%=================================================================
\section{Setup}
%=================================================================

\begin{defbox}
\begin{definition}[Even Integers]
Let $E$ denote the set of \textbf{even integers}:
\[
  E \;=\; \bigl\{\, z \in \Z \;\big|\; (\exists\, k \in \Z)\;(2k = z) \,\bigr\}
    \;=\; \{\dots,\, -4,\, -2,\, 0,\, 2,\, 4,\, \dots\}.
\]
\end{definition}
\end{defbox}

\begin{notebox}
\begin{remark}
We take $\N = \{0, 1, 2, \ldots\}$.  The goal is to prove $|\N| = |E|$ by
constructing an explicit bijection $f\colon \N \to E$.  The key idea is to
\emph{interleave}: map even-indexed naturals to the non-negative even integers
($0, 2, 4, \ldots$) and odd-indexed naturals to the negative even integers
($-2, -4, -6, \ldots$).
\end{remark}
\end{notebox}

%=================================================================
\section{The Theorem}
%=================================================================

\begin{thmbox}
\begin{theorem}
$|\N| = |E|$.
\end{theorem}
\begin{proof}
Define $f\colon \N \to E$ by
\[
  f(n) \;=\;
  \begin{cases}
    n   & \text{if } 2 \mid n, \\
    -n-1 & \text{if } 2 \nmid n.
  \end{cases}
\]

\medskip
\noindent\textbf{Well-defined (maps into $E$).}
\begin{itemize}[leftmargin=2em, itemsep=2pt]
  \item If $2 \mid n$, then $f(n) = n$ is even, so $f(n) \in E$.
  \item If $2 \nmid n$, then $n$ is odd, so $n + 1$ is even, hence
        $f(n) = -(n+1)$ is even, so $f(n) \in E$.
\end{itemize}

\medskip
\noindent\textbf{Injective.}
We must show $(\forall\, a, b \in \N)\bigl(f(a) = f(b) \Rightarrow a = b\bigr)$.
Let $a, b \in \N$ and assume $f(a) = f(b)$.  We split on the sign of $f(a)$.

\medskip
\noindent\underline{Case 1: $f(a) \geq 0$.}

First we show $2 \mid a$ (i.e.\ $a$ is even).
Suppose, towards a contradiction, that $2 \nmid a$ (i.e.\ $a$ is odd).
Then by definition of $f$, $f(a) = -a - 1 \leq -1 < 0$ since $a \in \N$,
contradicting $f(a) \geq 0$.  Therefore $2 \mid a$.

Since $f(a) = f(b) \geq 0$, by a similar argument $2 \mid b$.

Since $2 \mid a$, by definition of $f$: $f(a) = a$.
Since $2 \mid b$, by definition of $f$: $f(b) = b$.
Therefore $a = f(a) = f(b) = b$, so $\boxed{a = b}$.

\medskip
\noindent\underline{Case 2: $f(a) < 0$.}

Suppose, towards a contradiction, that $2 \mid a$.
Then $f(a) = a \geq 0$ since $a \in \N$, contradicting $f(a) < 0$.
Therefore $2 \nmid a$, i.e.\ $2 \mid a - 1$.

Since $f(a) = f(b) < 0$, by a similar argument $2 \mid b - 1$.

By definition of $f$: $f(a) = -a - 1$ and $f(b) = -b - 1$.
So $-a - 1 = -b - 1$, which gives $-a = -b$, so $\boxed{a = b}$.

\medskip
In both cases $f(a) = f(b) \Rightarrow a = b$, so $f$ is injective.

\medskip
\noindent\textbf{Surjective.}
We must show $(\forall\, y \in E)(\exists\, x \in \N)(f(x) = y)$.
Let $y \in E$.  By definition of $E$, there exists $k \in \Z$ such that $2k = y$.
We split on the sign of $y$.

\medskip
\noindent\underline{Case 1: $y \geq 0$.}

Then $2k = y \geq 0$ means $k \geq 0$, so $k \in \N$.  Therefore $2k \in \N$.
Take $x = 2k$.  Since $2 \mid 2k$, by definition of $f$:
\[
  f(x) = f(2k) = 2k = y.
\]

\medskip
\noindent\underline{Case 2: $y < 0$.}

Since $y \in \Z$ and $y < 0$, we have $y \leq -1$.
Then $y + 1 \leq 0$, so $-(y+1) \geq 0$.
Since $-(y+1) \in \Z$ and $-(y+1) \geq 0$, we have $-(y+1) \in \N$.

Now observe: since $y = 2k$,
\[
  -(y+1) = -(2k+1) = -2k - 1.
\]
So
\[
  -(y+1) - 1 = -2k - 1 - 1 = -2k - 2 = 2(-k-1),
\]
which shows $2 \mid -(y+1) - 1$, i.e.\ $2 \nmid -(y+1)$ (since $-(y+1)$ is odd).
Therefore by definition of $f$, with $x = -(y+1)$:
\[
  f(x) = f(-(y+1)) = -(-(y+1)) - 1 = (y+1) - 1 = y.
\]

\medskip
In both cases we found $x \in \N$ with $f(x) = y$, so $f$ is surjective.

\medskip
Therefore $f$ is injective and $f$ is surjective, so $f$ is bijective by definition.
Hence $|E| = |\N|$ by definition of cardinality equality. $\qed$
\end{proof}
\end{thmbox}

%=================================================================
\section{Explicit Correspondence}
%=================================================================

\begin{notebox}
\begin{example}
The first several values of $f$:

\medskip
\begin{center}
\renewcommand{\arraystretch}{1.5}
\begin{tabular}{|c|c|c|c|c|c|c|c|c|c|}
\hline
$n$    & $0$ & $1$  & $2$ & $3$  & $4$ & $5$  & $6$ & $7$  & $\cdots$ \\
\hline
$f(n)$ & $0$ & $-2$ & $2$ & $-4$ & $4$ & $-6$ & $6$ & $-8$ & $\cdots$ \\
\hline
\end{tabular}
\end{center}

\medskip
Even naturals hit $0, 2, 4, \ldots$ and odd naturals hit $-2, -4, -6, \ldots$,
together covering every element of $E$ exactly once.
\end{example}
\end{notebox}

\end{document}